\documentclass[12pt,a4paper]{article}
\usepackage{tabularx}
\usepackage{booktabs}
\usepackage{longtable}
\usepackage{ltxtable}
\usepackage[latin1]{inputenc}
\usepackage{amssymb}
\usepackage[]{graphicx,rotating}
%\usepackage[T1]{fontenc}
\usepackage{parskip}
\usepackage{natbib}
\usepackage[official]{eurosym}
%\usepackage{mathrsfs}
%\usepackage{amsmath}
\usepackage{verbatim}
\usepackage{epstopdf} % to include eps files
\usepackage{enumitem} % for the abbreviations table
\usepackage{helvet} %Arial
\renewcommand{\familydefault}{\sfdefault} %Arial
%\usepackage{mathptmx} %Times New Roman
%\renewcommand{\familydefault}{cmr} % Computer Modern. alternative: cmss

\usepackage[left=2.5cm, right=2.5cm, top=2.5cm]{geometry}

\pagestyle{empty}
\parindent 0.0cm
\renewcommand{\baselinestretch}{1}
\newcommand{\bs}{\boldsymbol}
\pdfminorversion=7
\bibliographystyle{agsm}

%nice quotes
\let\oldquote\quote
\let\endoldquote\endquote
\renewenvironment{quote}[2][]
{\if\relax\detokenize{#1}\relax
	\def\quoteauthor{#2}%
	\else
	\def\quoteauthor{#2~---~#1}%
	\fi
	\oldquote}
{\par\nobreak\smallskip\hfill(\quoteauthor)%
	\endoldquote\addvspace{\bigskipamount}}

\begin{document}

\begin{center}
% \vspace*{1cm}
 \includegraphics[width=0.2\textwidth]{external/GU-Logo-blau-CMYK.eps} \vspace{1.5cm}
  
{\large{\bf Economic Consequences of the 1918 Flu Pandemic:\\
		Reviewing Evidence from the Historic Precedent of COVID-19}}

\textbf{Lukas J\"urgensmeier} \\
{\footnotesize (Mat.-Nr.: 6904281)}

  Term Paper \\\vspace{1.5cm}
  
  \begin{abstract}
  	The ongoing coverage of Covid-19 often describes the current crisis as \textit{unprecedented}.
  	However, there exists precedent: The 1918 flu pandemic.
  	This paper first examines the economic consequences of the pandemic one century ago.
  	Second, it highlights to which extent lessons from the 1918 flu pandemic could be transferred to efficiently handle the current crisis.
  \end{abstract} \vspace{1.5cm}
  
  submitted to \\\vspace{0.5cm}
  \textbf{Prof. Dr. Volker Casperi} \\
  Graduate School of Economics, Finance, and Management \\
  Goethe University \\
  Frankfurt am Main \vspace{1.5cm}
  
  in fulfillment of the requirements \\
  of the lecture \\\vspace{0.5cm}
  \textbf{Historical and Normative Foundations of Economics} \\\vspace{0.5cm}
  July 30, 2020
  
\end{center}


\pagebreak
\pagestyle{plain}
\pagenumbering{Roman}
\tableofcontents
\pagebreak
\listoffigures
\listoftables
\newpage
\newlist{abbrv}{itemize}{1}
\setlist[abbrv,1]{label=,labelwidth=1in,align=parleft,itemsep=0.1\baselineskip,leftmargin=!}
 
\section*{List of Abbreviations}
%\chaptermark{List of Abbreviations}
 
\begin{abbrv}
 
\item[ABC]			Placeholder

 
\end{abbrv}
\newpage
\setcounter{page}{2}
\pagenumbering{arabic}
\setlength{\baselineskip}{1.5\baselineskip}
\pagestyle{plain}


\section{Introduction}

\begin{quote}{Mark Twain, 1910}
	\textit{"History doesn't repeat itself, but it rhymes."}
\end{quote}


\section{Development of the 1918 Flu Pandemic}

The late 1910s were marked by devastation.
The first truly global conflict, World War I, exerted a huge human toll between 1914 and 1918:
Historians estimate approximately 40 million civilian and military causalities resulting from the war \citep{royde-smithWorldWarKilled2020}.
However, a significant share of those deaths is not directly attributed to military conflict,
but to the second devastating event overlapping with the end of the military conflict and lasting at least until 1919: The 1918 flu pandemic.

Beginning in March 1918, a novel influenza virus emerged in mid-western parts of the United States and spread virulently, possibly sped up by large troop movements to Europe due to World War I. It reached almost every part of the world within six months \citep{pattersonGeographyMortality19181991}.
In mainly three waves in spring 1918, fall 1918, and winter 1918/19 in the northern hemisphere, the virus infected approximately 500 million people, roughly one third of the entire world population at that time \citep{taubenberger1918InfluenzaMother2006}.
Earlier research estimated the global death toll from the virus between 21.5 \cite{jordanEpidemicInfluenzaSurvey1927} and later up to 39.3 million \citep{pattersonGeographyMortality19181991}.
More recent studies, while acknowledging the large uncertainty, account for greater under-reporting of causalities and hence estimate the global death toll between 50 and 100 million \citep{johnsonUpdatingAccountsGlobal2002}.
It is therefore reasonable to assume that the health emergency caused more causalities that the enormous toll of World War I.

While influenza outbreaks before and after the 1918 flu pandemic mainly caused worse courses of the disease and a larger death toll on the very young and older population ("U-shape"),
the outbreak in 1918 displayed an uncommon W-shape by additionally killing many otherwise healthy young adults aged 20-40 \citep{taubenberger1918InfluenzaMother2006}\footnote{
	see figures 2 and 3 in \cite{taubenberger1918InfluenzaMother2006} for detailed age profiles of incidences and causalities of the 1918 flu pandemic in comparison to the normal influenza shortly before and after.}.

Especially this ruthlessness on people in their prime working age came with significant implications for the economy.
While this section focused on the broader health consequences of the pandemic, the next part examines the economic consequences of the devastating 1918 flu pandemic.

\section{Economic Consequences of the 1918 Flu Pandemic}
\subsection{Short-term: During the Outbreak in 1918-19}
\subsubsection{Health Care Costs}
\subsection{Absenteeism}

\begin{itemize}
	\item Huge effect on absentee rates. Especially noteworthy in industries deemed essential such as shipyards. According to \cite{barryPandemicsAvoidingMistakes2009}, "absentee rates in the shipyards -- one of the few industries for which there are good data -- still ranged from 45\% to 58\%" \citep{turnerReportPreventiveMeasure}.
	Those high rates were reached even though the incentives were aligned such that they had all economic reasons to report to work: Employees only received their wages if they worked.
\end{itemize}

\subsubsection{Non-Pharmaceutical Interventions}

\begin{itemize}
	\item Debate about effectiveness of NPIs on health outcomes:
	\cite{markelNonpharmaceuticalInterventionsImplemented2007}: State that US cities with early and continuous NPI measures such as quarantines overall had better health outcomes.
	\cite{barryCommentsNonpharmaceuticalInterventions2007}, one of the leading scholars on the 1918 flu pandemic, in a direct reply, however questions the methodology of the study.
	While he does not per se challenge the notion that NPIs did have a positive effect on slowing the spread and decreasing the negative health outcomes of the pandemic, he concludes that their methodology does not hold up to the burden of proof.
	This debate shows how difficult a causal statement about the usefulness of NPIs is in this case due to limited historical data.
	
	\item \cite{garrettEconomicEffects19182007}: Reviews newspaper articles that cite 50\% decrease in revenue, unemployment, incapacitated workers due to illness
	
	\item \cite{correiaPandemicsDepressEconomy2020}:
	Using geographic variation in mortality during the 1918 Flu Pandemic in
	the U.S., we find that more exposed areas experience a sharp and persistent decline
	in economic activity. The estimates imply that the pandemic reduced manufacturing
	output by 18 percent. The downturn is driven by both supply and demand-side channels.
	Further, building on findings from the epidemiology literature establishing that NPIs
	decrease influenza mortality, we use variation in the timing and intensity of NPIs
	across U.S. cities to study their economic effects. We find that cities that intervened
	earlier and more aggressively do not perform worse and, if anything, grow faster after
	the pandemic is over. Our findings thus indicate that NPIs not only lower mortality;
	they also mitigate the adverse economic consequences of a pandemic
	
	Altogether, our	findings suggest that pandemics can have substantial economic costs, and NPIs can have
	economic merits, beyond lowering mortality.
\end{itemize}

\subsection{Medium- to Long-term: After the Immediate Health Thread Subsided}

\subsubsection{Negative Impact on Fetal Health}
\begin{itemize}
	\item \cite{almond1918InfluenzaPandemic2006}: cohorts in utero during
	the pandemic displayed reduced educational attainment, increased
	rates of physical disability, lower income, lower socioeconomic status,
	and higher transfer payments compared with other birth cohorts.
	
	likelihood of being poor rose as much as 15 percent
	
	Prenatal exposure to the 1918 influenza pandemic had large negative
	effects on adult economic outcomes. 
\end{itemize}

\subsubsection{Loss of Life}

\begin{itemize}
	\item \cite{fanInclusiveCostPandemic2016}: In this paper we use findings from that literature to generate an
	estimate of pandemic cost that is inclusive of both income loss and the cost of elevated mortality.
	
	For moderately severe pandemics about 40\% of inclusive cost results from income loss. For
	severe pandemics this fraction declines to 12\%: the intrinsic cost of elevated mortality becomes
	completely dominant.
\end{itemize}

\subsubsection{Interest Rates, Wages, and Asset Prices}

\begin{itemize}
	\item \cite{correiaPandemicsDepressEconomy2020} uses variation in timing of NPIs to estimate economic effect of the pandemic. Cities with strong NPI intervention showed faster recovery and less adverse effects on the economy than comparable areas with less strong NPIs.
	\item \cite{brainerdEconomicEffects19182003}: large and robust positive effect of the influenza epidemic on per capita income growth across states during the 1920s. 
	\item \cite{karlssonImpact1918Spanish2014}: We find that the pandemic led to a significant increase in poverty rates.
	There is also relatively strong evidence that capital returns were negatively affected by the pandemic.
	 However, we find robust evidence that the influenza had no discernible effect on earnings.
	 Thisfinding is surprising since it goes against most previous empirical studies as well as theoretical predictions.
	
	\item \cite{jordaLongerrunEconomicConsequences2020}: We study rates of return on assets
	using a dataset stretching back to the 14th century, focusing on 15 major pandemics
	where more than 100,000 people died. In addition, we include major armed conflicts
	resulting in a similarly large death toll. Significant macroeconomic after-effects of the
	pandemics persist for about 40 years, with real rates of return substantially depressed.
	In contrast, we find that wars have no such effect, indeed the opposite. This is consistent
	with the destruction of capital that happens in wars, but not in pandemics. Using
	more sparse data, we find real wages somewhat elevated following pandemics. The
	findings are consistent with pandemics inducing labor scarcity and/or a shift to greater
	precautionary savings.
\end{itemize}

\section{Parallels and Differences of the 1918 Flu Pandemic and Covid-19}
\begin{itemize}
	\item Contrary to the 1918 flu pandemic, COVID-19 mostly affects the older population that is about to drop out or has already left the labor force.
	Hence, there might be no sustained negative impact on labor supply and thus no positive effect on wages.
	\item Insurance modelers developed a framework to simulate an outbreak of the 1918 flu pandemic today \citep{madhavModelingModernDay2013}. Differences in population groups affected. 1918 mostly young adults and few older, 2020 mostly older adults.	Implications for different economic consequences: Today, less adults in the phase of prime economic productivity seem die from the disease.
	\item Timing right after World War II: Economy was already in a recession, which masked the economic shock due to the pandemic.
	Contrary, Covid-19 hit the world economy during a decade-long expansion.
	Hence, the economic stakes were higher and the immediate impact on the economy more visible.
\end{itemize}






\section{Conclusion}

\clearpage
\appendix
\section{Appendix}

\clearpage
\bibliography{library}

\newpage
\thispagestyle{empty}
\section*{Statutory Declaration}

Ich versichere hiermit, dass ich die vorliegende Arbeit selbs\"andig und ohne Benutzung anderer als der angegebenen Quellen und Hilfsmittel verfasst habe. W\"ortlich \"ubernommene S\"atze oder Satzteile sind als Zitat belegt, andere Anlehnungen, hinsichtlich Aussage und Umfang, unter Quellenangabe kenntlich gemacht. Die Arbeit hat in gleicher oder \"ahnlicher Form noch keiner Pr\"ufungsbeh\"orde vorgelegen und ist nicht ver\"offentlicht. Sie wurde nicht, auch nicht auszugsweise, f\"ur eine andere Pr\"ufungs- oder Studienleistung verwendet.


I herewith declare that I have completed the present term paper independently, without making use of
other than the specified literature and aids. Sentences or parts of sentences quoted literally are
marked as quotations; identification of other references with regard to the statement and scope of
the work is quoted. The thesis in this form or in any other form has not been submitted to an examination body and has not been published.
This thesis has not been used, either in whole or part, for another examination achievement.

\vspace{1cm}

Frankfurt am Main, July 30, 2019

\includegraphics[scale=0.12]{external/signature.png}\\
Lukas J\"urgensmeier
\end{document}
